%% ===================================================================
\section{Framework}\label{sec:framework}
%% ===================================================================

\subsection{Soft gluons, valence charges and the light-cone Hamiltonian}\label{sec:hamiltonian}

We work in light-cone gauge $A^+=0$ with the projectile moving in the $+$ direction. The gluon field is split into soft modes, with longitudinal momenta in the window $\Lambda<k^+<\vee$, and faster (valence) modes with $k^+>\vee$~\cite{Kovner:2005nq,Lublinsky:2016meo}:
\begin{equation}
A_i^a(x)=\int_{\Lambda}^{\vee}\frac{dk^+}{2\pi}\int\frac{d^2\mathbf k}{(2\pi)^2}\frac{1}{\sqrt{2k^+}}\Big(a_i^a(k^+,\mathbf k)\,e^{-ik\cdot x}+a_i^{a\dagger}(k^+,\mathbf k)\,e^{ik\cdot x}\Big)+\big(\text{modes with }k^+>\vee\big).
\label{eq:Asoft}
\end{equation}
In the eikonal (CGC) approximation the valence modes enter only through their colour charge density $\rho^a(x)$, a function of the transverse position $x$. Its components obey the $SU(N_c)$ algebra
\begin{equation}
[\rho^a(x),\rho^b(y)]=if^{abc}\rho^c(x)\,\delta^{(2)}(x-y).
\label{eq:rhoalg}
\end{equation}
The soft creation and annihilation operators satisfy $[a^a_i(k),a^{b\dagger}_j(p)]=(2\pi)^3\delta^{ab}\delta_{ij}\delta^{(3)}(k-p)$, with $\delta^{(3)}(k-p)=\delta(k^+-p^+)\delta^{(2)}(\mathbf k-\mathbf p)$. Throughout, $x,y,z,u,v,w,\ldots$ denote two-dimensional transverse positions, $\int_x\equiv\int d^2x$, and the transverse Fourier conventions are
\begin{equation}
a(k^+,\mathbf k)=\int d^2z\,e^{-i\mathbf k\cdot z}a(k^+,z),\qquad a^\dagger(k^+,\mathbf k)=\int d^2z\,e^{i\mathbf k\cdot z}a^\dagger(k^+,z),\qquad f(\mathbf k)=\int d^2x\,e^{i\mathbf k\cdot x}f(x),
\label{eq:FTconv}
\end{equation}
so that $[a^a_i(k^+,z),a^{b\dagger}_j(p^+,z')]=2\pi\,\delta^{ab}\delta_{ij}\delta(k^+-p^+)\delta^{(2)}(z-z')$.

The light-cone Hamiltonian is $H=H_0+H_{\rm int}$. The free part gives a soft gluon the energy $E_g(k)=\mathbf k^2/2k^+$. The interaction terms that contribute at the order we work are~\cite{Lublinsky:2016meo}:
\begin{itemize}
\item $H_g$, the emission and absorption of a soft gluon by the valence charges,
\begin{equation}
H_g=\int_{\Lambda}^{\vee}\frac{dk^+}{2\pi}\int\frac{d^2\mathbf k}{(2\pi)^2}\,\frac{g\,\mathbf k^i}{\sqrt2\,|k^+|^{3/2}}\Big[a_i^{a\dagger}(k^+,\mathbf k)\,\rho^a(-\mathbf k)+a_i^a(k^+,\mathbf k)\,\rho^a(\mathbf k)\Big];
\label{eq:Hg}
\end{equation}
\item $H_{ggg}$, the three-gluon vertex among soft gluons;
\item $H_{gg\text{-inst}}$, the instantaneous interaction of two soft gluons with the valence charge (from the longitudinal Coulomb field in light-cone gauge).
\end{itemize}
Their matrix elements between Fock states are collected in App.~\ref{app:conv}, Eqs.~\eqref{eq:m1}--\eqref{eq:m5}. The quartic soft vertex and quark loops are not needed for gluon production at this order, except through the gluon self-energy, where we keep only the gluon loop ($n_f=0$).

\subsection{The dressed projectile and the evolution operator \texorpdfstring{$\Omega$}{Omega}}\label{sec:Omega}

Let $|0\rangle$ be the state with the valence charges and no soft gluons. The physical projectile is the eigenstate of $H$ that reduces to $|0\rangle$ at $g=0$. Light-cone perturbation theory~\cite{Lepage:1980fj,Brodsky:1997de} constructs it order by order:
\begin{equation}
|\Psi\rangle=\mf N|0\rangle+\sum_n|n\rangle\frac{\langle n|H_{\rm int}|0\rangle}{E_0-E_n}+\sum_{n,m}|n\rangle\frac{\langle n|H_{\rm int}|m\rangle\langle m|H_{\rm int}|0\rangle}{(E_0-E_n)(E_0-E_m)}+\dots,\qquad E_0=0,
\label{eq:LCPT}
\end{equation}
where $|n\rangle,|m\rangle$ are free Fock states of soft gluons and $\mf N$ is fixed by normalization. The same construction dresses any soft state, e.g.\ a single soft gluon $|g^a_i(k)\rangle$. We write all dressed states with one unitary operator $\Omega$,
\begin{equation}
|\Psi\rangle=\Omega|0\rangle,\qquad |\Psi^a_i(k)\rangle=\Omega\,|g^a_i(k)\rangle,\qquad \Omega^\dagger\Omega=1.
\label{eq:Omegadef}
\end{equation}
$\Omega$ acts on the soft Fock space and depends on the valence charges $\rho$. Normal ordering the soft operators gives~\cite{Lublinsky:2016meo}
\begin{align}
\Omega={}&\mf N+\int\frac{d^3p}{(2\pi)^3}\Big[\mf A^b_{1j}(p)\,a^b_j(p)+\mf A^b_{2j}(p)\,a^{b\dagger}_j(p)\Big]\notag\\
&+\int\frac{d^3p}{(2\pi)^3}\frac{d^3q}{(2\pi)^3}\Big[\mf B^{cb}_{1kj}(p,q)\,a^b_j(p)a^c_k(q)+\mf B^{cb}_{2kj}(p,q)\,a^{b\dagger}_j(p)a^{c\dagger}_k(q)+\mf B^{cb}_{3kj}(p,q)\,a^{c\dagger}_k(q)a^b_j(p)\Big]\notag\\
&+\int\frac{d^3p}{(2\pi)^3}\frac{d^3q}{(2\pi)^3}\Big[\mf C^{dcb}_{1lkj}(p,q)\,a^{d\dagger}_l(q)a^{c\dagger}_k(p-q)a^b_j(p)+\mf C^{dcb}_{2lkj}(p,q)\,a^{b\dagger}_j(p)a^c_k(p-q)a^d_l(q)\Big]+\dots
\label{eq:OmegaExp}
\end{align}
The coefficients are operators in the valence colour space. Bose symmetry implies $\mf B^{cb}_{1kj}(p,q)=\mf B^{bc}_{1jk}(q,p)$, $\mf B^{cb}_{2kj}(p,q)=\mf B^{bc}_{2jk}(q,p)$ and $\mf C^{dcb}_{1lkj}(p,q)=\mf C^{cdb}_{1klj}(p,p-q)$ (and the same for $\mf C_2$). The dots stand for terms with more soft operators, which do not contribute to single gluon production at order $g^4$. The coefficients mix orders in $g$. We need $\mf N=1+\mf N_2$ to order $g^2$ and
\begin{equation}
\mf A_{1,2}=\mf A^{(1)}_{1,2}+\mf A^{(3)}_{1,2}
\end{equation}
to order $g^3$. $\mf B_{1,2,3}$ are needed at order $g^2$ and $\mf C_{1,2}$ at order $g$.

\paragraph*{Unitarity.} The condition $\Omega^\dagger\Omega=1$, normal ordered, relates the daggered and undaggered coefficients. To the order needed, and with $\mf N=\mf N^\dagger$,
\begin{align}
&\mf N=1-\frac12\int\frac{d^3p}{(2\pi)^3}\,\mf A^{b}_{1j}(p)\,\mf A^{\dagger b}_{1j}(p),\label{eq:unitN}\\
&\mf A^{(1)\dagger b}_{1j}(p)=-\mf A^{(1)b}_{2j}(p),\label{eq:unitA1}\\
&\mf A^{(3)\dagger b}_{1j}(p)=-\mf A^{(3)b}_{2j}(p)-\big[\mf N_2,\mf A^{(1)b}_{2j}(p)\big]+\int\frac{d^3q}{(2\pi)^3}\Big[2\,\mf A^{(1)c}_{1k}(q)\,\mf B^{cb}_{2kj}(p,q)-\mf B^{\dagger cb}_{3kj}(p,q)\,\mf A^{(1)c}_{2k}(q)\notag\\
&\hspace{20em}-2\,\mf C^{\dagger dcb}_{1lkj}(p,q)\,\mf B^{cd}_{2kl}(q,p-q)\Big],\label{eq:unitA3}\\
&\mf B^{\dagger cb}_{1kj}(p,q)=\frac12\Big[-2\,\mf B^{cb}_{2kj}(p,q)+\mf A^{(1)b}_{2j}(p)\,\mf A^{(1)c}_{2k}(q)+\mf A^{(1)c}_{2k}(q)\,\mf A^{(1)b}_{2j}(p)+2\,\mf C^{cbd}_{1kjl}(q+p,q)\,\mf A^{(1)d}_{2l}(q+p)\Big],\label{eq:unitB1}\\
&\mf C^{\dagger dcb}_{2lkj}(p,q)=-\mf C^{dcb}_{1lkj}(p,q).\label{eq:unitC2}
\end{align}
The commutator in Eq.~\eqref{eq:unitA3} does not vanish. $\mf N_2\propto\int K\,\rho^a\rho^a$ [Eq.~\eqref{eq:NX} below] commutes with the total charge $\int_x\rho^a(x)$, but not with the local charge $\rho^b(u)$ in $\mf A^{(1)}$:
\begin{equation}
\Big[\int_{x,y}K(x,y)\,\rho^a(x)\rho^a(y),\ \rho^b(u)\Big]=if^{abc}\int_xK(x,u)\,\big\{\rho^a(x),\rho^c(u)\big\}
\label{eq:Ncomm}
\end{equation}
for a symmetric kernel $K(x,y)=K(y,x)$. It is a two-charge term proportional to $\log(\vee/\Lambda)$ [the logarithm of $\mf N_2$], and therefore contributes only to the leading-logarithmic part of the cross section, never to the finite part of Sec.~\ref{sec:finite}. In Sec.~\ref{sec:LL} the LL part is derived directly from the leading-log form of $\Omega$, in which all such orderings are automatically included.
\note{The draft notes set $[\mf N_2,\mf A^{(1)}]=0$; with Eq.~\eqref{eq:Ncomm} the Group~I bookkeeping of the $\log(\vee/\Lambda)$ terms has to include this commutator.}
%TODO: propagate [N_2, A^(1)] into the row-by-row LL bookkeeping of Group I (4rho notes).
The coefficients $\mf A^{(1)}_1$, $\mf A^{(3)}_2$, $\mf B_2$, $\mf B_3$ and $\mf C_1$ are obtained from LCPT; $\mf A^{(3)}_1$, $\mf B_1$ and $\mf C_2$ then follow from Eqs.~\eqref{eq:unitA3}--\eqref{eq:unitC2}. The product $\mf A^{(1)}\mf A^{(1)}$ in Eq.~\eqref{eq:unitB1} is kept as an operator product: its two orderings differ by the commutator~\eqref{eq:rhoalg}, which is of the same order as the $\rho$-linear part of $\mf B_2$.

\paragraph*{Coordinate space.} With the conventions~\eqref{eq:FTconv}, the expansion~\eqref{eq:OmegaExp} becomes
\begin{align}
\Omega={}&\mf N+\int\frac{dp^+}{2\pi}\int_x\Big[\mf A^{(1)b}_{1j}(p^+,x)\,a^b_j(p^+,x)-\mf A^{\dagger(1)b}_{1j}(p^+,x)\,a^{\dagger b}_j(p^+,x)\notag\\
&\hspace{7em}+\mf A^{(3)b}_{1j}(p^+,x)\,a^b_j(p^+,x)+\mf A^{(3)b}_{2j}(p^+,-x)\,a^{\dagger b}_j(p^+,x)\Big]\notag\\
&+\int\frac{dp^+}{2\pi}\frac{dq^+}{2\pi}\int_{x,y}\Big[\mf B^{cb}_{1kj}(p^+,x;q^+,y)\,a^b_j(p^+,x)a^c_k(q^+,y)\notag\\
&\hspace{7em}+\mf B^{cb}_{2kj}(p^+,-x;q^+,-y)\,a^{\dagger b}_j(p^+,x)a^{\dagger c}_k(q^+,y)\notag\\
&\hspace{7em}+\mf B^{cb}_{3kj}(p^+,x;q^+,-y)\,a^{\dagger c}_k(q^+,y)a^b_j(p^+,x)\Big]\notag\\
&+\int\frac{dp^+}{2\pi}\frac{dq^+}{2\pi}\int_{x,y,z}\Big[\mf C^{dcb}_{1lkj}(p^+,z-y;q^+,y-x)\,a^{\dagger d}_l(q^+,x)a^{\dagger c}_k(p^+-q^+,y)a^b_j(p^+,z)\notag\\
&\hspace{7em}-\mf C^{\dagger dcb}_{1lkj}(p^+,z-y;q^+,y-x)\,a^{\dagger b}_j(p^+,z)a^c_k(p^+-q^+,y)a^d_l(q^+,x)\Big]+\dots,
\label{eq:OmegaX}
\end{align}
where we have used Eqs.~\eqref{eq:unitA1} and~\eqref{eq:unitC2}. The unitarity relations in coordinate space read
\begin{align}
\mf A^{(1)\dagger b}_{1j}(p^+,z)={}&-\mf A^{(1)b}_{2j}(p^+,-z),\label{eq:unitA1X}\\
\mf A^{(3)\dagger b}_{1j}(p^+,x)={}&-\mf A^{(3)b}_{2j}(p^+,-x)-\big[\mf N_2,\mf A^{(1)b}_{2j}(p^+,-x)\big]\notag\\
&+\int\frac{dq^+}{2\pi}\int_y\Big[2\,\mf A^{(1)c}_{1k}(q^+,y)\,\mf B^{cb}_{2kj}(p^+,-x;q^+,-y)\notag\\
&\hspace{5em}-\mf B^{\dagger cb}_{3kj}(p^+,x;q^+,-y)\,\mf A^{(1)c}_{2k}(q^+,-y)\notag\\
&\hspace{5em}-2\int_z\mf C^{\dagger dcb}_{1lkj}(p^+,x-z;q^+,z-y)\,\mf B^{dc}_{2lk}(p^+-q^+,-z;q^+,-y)\Big],\label{eq:unitA3X}\\
\mf B^{\dagger cb}_{1kj}(p^+,x;q^+,y)={}&\frac12\Big[-2\,\mf B^{cb}_{2kj}(p^+,-x;q^+,-y)+\mf A^{(1)b}_{2j}(p^+,-x)\,\mf A^{(1)c}_{2k}(q^+,-y)\notag\\
&\hspace{3em}+\mf A^{(1)c}_{2k}(q^+,-y)\,\mf A^{(1)b}_{2j}(p^+,-x)\notag\\
&\hspace{3em}+2\int_z\mf C^{cbd}_{1kjl}(q^++p^+,z-x;q^+,x-y)\,\mf A^{(1)d}_{2l}(q^++p^+,-z)\Big],\label{eq:unitB1X}\\
\mf C^{\dagger dcb}_{2lkj}(p^+,z;q^+,x)={}&-\mf C^{dcb}_{1lkj}(p^+,-z;q^+,-x).\label{eq:unitC2X}
\end{align}

\paragraph*{Action of $\Omega$ on the vacuum and on one gluon.} The coefficients are read off from
\begin{align}
\Omega|0\rangle={}&\mf N|0\rangle+\frac{1}{(2\pi)^{1/2}}\int dp^+\int_x\Big[\mf A^{(3)b}_{2j}(p^+,-x)-\mf A^{\dagger(1)b}_{1j}(p^+,x)\Big]|g^b_j(p^+,x)\rangle\notag\\
&+\frac{1}{2\pi}\int dp^+dq^+\int_{z,x}\mf B^{cb}_{2kj}(p^+,-z;q^+,-x)\,|g^b_j(p^+,z)\,g^c_k(q^+,x)\rangle,\label{eq:Omegavac}\\[3pt]
\Omega|g^a_i(k^+,z)\rangle={}&\frac{1}{(2\pi)^{1/2}}\Big[\mf A^{(1)a}_{1i}(k^+,z)+\mf A^{(3)a}_{1i}(k^+,z)\Big]|0\rangle+\mf N|g^a_i(k^+,z)\rangle\notag\\
&+\frac{1}{2\pi}\int dp^+\int_x\mf B^{ba}_{3ji}(k^+,z;p^+,-x)\,|g^b_j(p^+,x)\rangle\notag\\
&+\frac{1}{(2\pi)^{1/2}}\int dp^+\int_x\Big[\mf A^{(3)b}_{2j}(p^+,-x)-\mf A^{\dagger(1)b}_{1j}(p^+,x)\Big]|g^a_i(k^+,z)\,g^b_j(p^+,x)\rangle\notag\\
&+\frac{1}{(2\pi)^{1/2}}\int dp^+\int_{x,y}\mf C^{bca}_{1jki}(k^+,z-y;p^+,y-x)\,|g^b_j(p^+,x)\,g^c_k(k^+-p^+,y)\rangle+\dots,\label{eq:Omegag}
\end{align}
with the states of App.~\ref{app:conv}. Comparing Eq.~\eqref{eq:Omegavac} with the LCPT vacuum wave function fixes $\mf N$, $\mf A^{(1)}$, $\mf A^{(3)}_2$ and $\mf B_2$. Comparing Eq.~\eqref{eq:Omegag} with the dressed one-gluon state fixes $\mf B_3$ and $\mf C_1$.

\subsection{Scattering and the inclusive cross section}\label{sec:xsecdef}

The target acts on the projectile through the eikonal $S$-matrix. Every colour charge, valence or soft, is rotated by the adjoint Wilson line of the target at its transverse position~\cite{Kovner:2005nq}:
\begin{equation}
S^\dagger\rho^a(x)S=U^{ab}(x)\rho^b(x),\quad S^\dagger a^a_i(k^+,x)S=U^{ab}(x)\,a^b_i(k^+,x),\quad S^\dagger a^{a\dagger}_i(k^+,x)S=U^{ab}(x)\,a^{b\dagger}_i(k^+,x).
\label{eq:Srot}
\end{equation}
Here $U(x)=\mathcal P\exp\{ig\int dx^+T^a\alpha^a(x^+,x)\}$, with $(T^a)_{bc}=-if^{abc}$ and $\alpha^a$ the target field. $U$ is real and orthogonal, $U^\dagger=U^T$. We denote a rotated operator by a bar, e.g.\ $\bar\Omega\equiv S^\dagger\Omega S$, which is $\Omega$ with $\rho\to U\rho$. Rotated coefficients are written $\bar{\mf A}$, $\bar{\mf B}$, $\bar{\mf C}$; for instance $\bar{\mf A}^{(1)}$ is $\mf A^{(1)}$ with $\rho^b(x)\to U^{bc}(x)\rho^c(x)$.

The incoming state is $\Omega|0\rangle$ and the outgoing state is $S\Omega|0\rangle$. The observable is the number of \emph{dressed} gluons in the outgoing state. The number operator of dressed gluons is $\Omega\,a^\dagger a\,\Omega^\dagger$, so
\begin{equation}
\frac{d^3N}{d^3k}=\frac{1}{(2\pi)^3}\int_{w,w'}e^{-ik\cdot(w'-w)}\,\langle0|\,\Omega^\dagger S^\dagger\Omega\;a^{a\dagger}_i(k^+,w)\,a^a_i(k^+,w')\;\Omega^\dagger S\,\Omega\,|0\rangle,
\label{eq:xsec}
\end{equation}
where $d^3k\equiv dk^+d^2\mathbf k$ and the expectation value includes the average over the valence charges and over the target fields. Inserting $SS^\dagger=1$ and complete sets of one- and two-gluon states gives
\begin{equation}
\frac{d^3N}{d^3k}=\frac{1}{(2\pi)^3}\bigg[\underbrace{\mathcal A^{\dagger a}_i(k)\,\mathcal A^a_i(k)}_{\text{one gluon}}+\int dq^+\int_v\underbrace{\mathcal A^{\dagger ab}_{ij}(k;q^+,v)\,\mathcal A^{ab}_{ij}(k;q^+,v)}_{\text{two gluons}}\bigg],
\label{eq:groups}
\end{equation}
with the amplitudes
\begin{align}
\mathcal A^a_i(k)&=\sqrt{2\pi}\int_{w'}e^{-ik\cdot w'}\,U^{ab}(w')\,\langle g^b_i(k^+,w')|\,\bar\Omega^\dagger\Omega\,|0\rangle,\label{eq:amp1}\\
\mathcal A^{ab}_{ij}(k;q^+,v)&=\sqrt{2\pi}\int_{w'}e^{-ik\cdot w'}\,U^{ac}(w')\,\langle g^b_j(q^+,v)\,g^c_i(k^+,w')|\,\bar\Omega^\dagger\Omega\,|0\rangle.\label{eq:amp2}
\end{align}
The operator $\bar\Omega^\dagger\Omega$ makes the structure of Eq.~\eqref{eq:xsec} transparent. $\Omega$ creates the soft cloud of the incoming projectile. The target rotates it, and $\bar\Omega^\dagger=S^\dagger\Omega^\dagger S$ removes the cloud that the \emph{rotated} valence charges carry in the final state. If the target is trivial, $U=1$, then $\bar\Omega^\dagger\Omega=1$ and nothing is produced. The second gluon in Eq.~\eqref{eq:amp2} is unobserved and is integrated over its momentum fraction $q^+$ and position $v$.

\subsection{Leading order}\label{sec:LO}

At order $g$ only $\mf A^{(1)}$ contributes. Using Eq.~\eqref{eq:unitA1X} and the coefficient~\eqref{eq:A1X} below,
\begin{align}
\mathcal A^a_i(k)&=\int_{w'}e^{-ik\cdot w'}\Big[-U^{ab}(w')\,\mf A^{\dagger(1)b}_{1i}(k^+,w')+\bar{\mf A}^{\dagger(1)a}_{1i}(k^+,w')\Big]\notag\\
&=\frac{ig}{\sqrt2\,\pi\sqrt{k^+}}\int_{w'}e^{-ik\cdot w'}\int_x\KK^i(x-w')\big[U(w')-U(x)\big]^{ab}\rho^b(x),
\label{eq:ampLO}
\end{align}
where we introduced the Weizs\"acker--Williams field
\begin{equation}
\KK^i(X)\equiv\frac{X^i}{X^2}.
\end{equation}
The two terms are the emission of the gluon before the target (rotated at $w'$) and after it (the emitter rotated at $x$). The LO spectrum is
\begin{align}
\frac{d^3N_{\rm LO}}{d^3k}={}&\frac{1}{(2\pi)^3}\,\frac{g^2}{2\pi^2k^+}\int_{w,w'}e^{-ik\cdot(w'-w)}\int_{x,x'}\KK(x-w)\cdot\KK(x'-w')\notag\\
&\times\Big\langle\rho^b(x)\Big[\big(\Ud(x)-\Ud(w)\big)\big(U(x')-U(w')\big)\Big]^{bb'}\rho^{b'}(x')\Big\rangle.
\label{eq:LOop}
\end{align}
For a projectile whose colour correlator is diagonal, $\langle\rho^b(x)\rho^{b'}(x')\rangle=\delta^{bb'}\langle\rho^f(x)\rho^f(x')\rangle/(N_c^2-1)$, this becomes the familiar form~\cite{Kovchegov:2001sc,Dumitru:2001ux,Blaizot:2004wu}
\begin{align}
\frac{d^3N_{\rm LO}}{d^3k}={}&\frac{\alpha_s}{4\pi^4(N_c^2-1)\,k^+}\int_{x,w,x',w'}e^{-ik\cdot(w'-w)}\,\KK(x-w)\cdot\KK(x'-w')\notag\\
&\times\langle\rho^f(x)\rho^f(x')\rangle\,\Tr\Big\{\big[\Ud(x)-\Ud(w)\big]\big[U(x')-U(w')\big]\Big\}.
\label{eq:LO}
\end{align}
Integrating over $\mathbf k$ sets $w=w'$ and gives $dN_{\rm LO}/dk^+=(\alpha_s/\pi^2)(1/k^+)\int_w\langle\dots\rangle$. Each additional emission of a gluon with momentum fraction $dq^+/q^+$ therefore costs $(\alpha_s/\pi^2)\,dq^+/q^+$, times the kernel $\int_z\KK(s-z)\cdot\KK(s'-z)$ for emitters at $s$ and $s'$. This normalization fixes the coefficients of the logarithms in Sec.~\ref{sec:LL}.

For the NLO terms it is convenient to write the LO spectrum in terms of the operator
\begin{equation}
\Ord_{\rm LO}\equiv\Big[U^{ab'}(x')\rho^{b'}(x')-U^{ab'}(w')\rho^{b'}(x')\Big]\Big[U^{ab}(w)\rho^b(x)-U^{ab}(x)\rho^b(x)\Big],\quad K_{\rm LO}\equiv\KK(x'-w')\cdot\KK(x-w),
\label{eq:OLO}
\end{equation}
in which the amplitude sits at $w$ and its conjugate at $w'$. In this labelling
\begin{equation}
\frac{d^3N_{\rm LO}}{d^3k}=-\pre\frac{g^2}{2\pi^2k^+}\int_{x,x'}K_{\rm LO}\,\Ord_{\rm LO},
\label{eq:LOrelab}
\end{equation}
which is positive since $-\Ord_{\rm LO}=\{[U(w')-U(x')]\rho(x')\}^a\{[U(w)-U(x)]\rho(x)\}^a$. Eq.~\eqref{eq:LOrelab} follows from Eq.~\eqref{eq:LOop} by $w\leftrightarrow w'$, $x\leftrightarrow x'$ and complex conjugation; the NLO terms below are always accompanied by their complex conjugates, so the two labellings are equivalent.
